% ═════════════════════════════════════════════════════════════════════════════
%  chapitres/chapitre3.tex — Threat Modeling et Analyse de Sécurité
% ═════════════════════════════════════════════════════════════════════════════
\chapter{Threat Modeling et Analyse de Sécurité}

\section*{Introduction}
\label{sec:introduction-ch3}

Ce chapitre applique la méthodologie STRIDE au prototype de système de paiement
déployé avec Istio, afin d'identifier et d'analyser les menaces de sécurité
spécifiques aux architectures microservices~[4,1]. Pour chacune des six catégories
de menaces, les contre-mesures apportées par le Service Mesh sont détaillées et
mises en correspondance avec les mécanismes Istio déployés.

% ─────────────────────────────────────────────────────────────────────────────
\section{Méthodologie STRIDE}
% ─────────────────────────────────────────────────────────────────────────────

Le \textbf{threat modeling} (modélisation des menaces) est une démarche structurée
d'identification et d'analyse des risques de sécurité d'un système. Parmi les
frameworks existants --- STRIDE, PASTA, LINDDUN, MITRE ATT\&CK --- nous avons
retenu STRIDE pour sa pertinence dans les architectures orientées services et sa
correspondance directe avec les mécanismes de protection qu'offre Istio~[4,1].

STRIDE est un acronyme développé par Microsoft qui catégorise les menaces en six
familles : \textbf{S}poofing (usurpation d'identité), \textbf{T}ampering
(altération de données), \textbf{R}epudiation (non-traçabilité des actions),
\textbf{I}nformation Disclosure (divulgation d'informations), \textbf{D}enial of
Service (déni de service) et \textbf{E}levation of Privilege (élévation de
privilèges). Chaque catégorie correspond à une propriété de sécurité à protéger.

% ─────────────────────────────────────────────────────────────────────────────
\section{Architecture du Système Analysé}
% ─────────────────────────────────────────────────────────────────────────────

Le système étudié est un prototype de système de paiement déployé sur
Kubernetes~[11,12] avec deux microservices communicants :

\begin{itemize}[topsep=4pt, itemsep=3pt]
  \item \textbf{payment-api} (Go, port 8080) : service principal exposé via
    gateway. Reçoit les requêtes de paiement, valide le JWT et transmet à
    \texttt{fraud-detection}.

  \item \textbf{fraud-detection} (Python/Flask, port 8080) : service interne de
    détection de fraude. Accessible uniquement depuis \texttt{payment-api} via
    SPIFFE~[7].
\end{itemize}

La surface d'attaque comprend trois périmètres distincts :

\begin{enumerate}[topsep=4pt, itemsep=3pt]
  \item Le trafic \textbf{externe} : client $\rightarrow$ gateway $\rightarrow$
    payment-api.
  \item Le trafic \textbf{interne} : payment-api $\rightarrow$ fraud-detection.
  \item Les tentatives d'accès depuis des \textbf{namespaces ou pods non
    autorisés}.
\end{enumerate}

% ─────────────────────────────────────────────────────────────────────────────
\section{Analyse des Six Catégories STRIDE}
% ─────────────────────────────────────────────────────────────────────────────

\subsection{Spoofing : Usurpation d'Identité}

Dans une architecture microservices sans Service Mesh, un service malveillant
peut se faire passer pour un service légitime sans aucun mécanisme de vérification
d'identité. Les appels HTTP entre services sont anonymes par nature.

Istio répond à cette menace via deux mécanismes complémentaires~[1,7] :

\begin{enumerate}[topsep=4pt, itemsep=3pt]
  \item Le \textbf{mTLS en mode STRICT} impose que chaque service présente un
    certificat valide lors de l'établissement de la connexion. Ces certificats
    sont générés automatiquement par \texttt{istiod} et encodent l'identité
    SPIFFE du workload sous la forme
    \texttt{spiffe://cluster.local/ns/\{namespace\}/sa/\{serviceaccount\}}~[1,7].

  \item Les \textbf{AuthorizationPolicies} filtrent les requêtes en fonction du
    principal SPIFFE source, garantissant que seul \texttt{payment-api} peut
    joindre \texttt{fraud-detection}.
\end{enumerate}

\subsection{Tampering : Altération de Données en Transit}

Sans chiffrement du trafic inter-services, un attaquant ayant accès au réseau
du cluster peut intercepter et modifier les données en transit. Dans le contexte
d'un système de paiement, cela pourrait permettre la modification des montants
de transactions ou des réponses de l'analyse de fraude.

Le mTLS garantit l'intégrité des données en transit via le protocole TLS~1.3.
Toute modification des données chiffrées invalide le \textit{Message
Authentication Code} (MAC) et provoque le rejet de la connexion. Par ailleurs,
la validation JWT assure l'intégrité des tokens d'authentification : un JWT dont
la signature a été modifiée est rejeté avec un code \texttt{401} par le composant
\texttt{RequestAuthentication} d'Istio.

\subsection{Repudiation : Non-Traçabilité des Actions}

La repudiation désigne la capacité d'un acteur malveillant à nier avoir effectué
une action. Dans les systèmes distribués, l'absence de logs centralisés et
horodatés facilite cette forme d'attaque.

Le proxy Envoy génère automatiquement un \textit{access log} détaillé pour chaque
requête traitée, incluant : l'horodatage précis, l'adresse IP source, l'URI, le
code de réponse HTTP, la durée, le principal SPIFFE source et la politique
d'autorisation appliquée. Ces logs, accessibles via \texttt{kubectl logs} et
visualisables dans Kiali, constituent un \textbf{audit trail} complet et non
manipulable par les applications~[8,10]. Aucune action dans le mesh ne peut être
effectuée sans laisser une trace dans les logs Envoy.

\subsection{Information Disclosure : Divulgation d'Information}

Sans chiffrement, le trafic inter-services est lisible par n'importe quel
processus ayant accès au réseau du cluster. Une capture de paquets
(\texttt{tcpdump}, Wireshark) révèle les données métier en clair : montants de
transactions, identifiants de cartes, résultats d'analyse de fraude.

Le mTLS STRICT imposé par la \texttt{PeerAuthentication} chiffre l'intégralité
du trafic inter-pods via TLS~1.3. Une capture réseau entre \texttt{payment-api}
et \texttt{fraud-detection} ne révèle que des \textit{TLS Application Data
Records} illisibles. Le certificat présenté lors du handshake TLS révèle
l'identité SPIFFE du service, mais aucune donnée applicative.

\subsection{Denial of Service : Déni de Service}

Les microservices exposés directement sur le réseau sont vulnérables aux attaques
par saturation. Un service inondé de requêtes peut devenir indisponible, et dans
une architecture en chaîne, cette indisponibilité peut se propager en cascade.

Istio apporte deux niveaux de protection complémentaires :

\begin{itemize}[topsep=4pt, itemsep=3pt]
  \item Les \textbf{DestinationRules} permettent de définir des
    \textit{connection pools} (limites de connexions simultanées et de requêtes
    en attente) au niveau du sidecar Envoy, sans dépendance à un service externe.

  \item Les \textbf{EnvoyFilters} permettent de configurer un mécanisme de
    \textit{rate limiting} par token bucket (ex. : 50~requêtes/min), bloquant
    les requêtes excédentaires avec un code \texttt{429}~[8].
\end{itemize}

Ces protections opèrent avant que la requête n'atteigne le code applicatif.

\subsection{Elevation of Privilege : Élévation de Privilèges}

L'élévation de privilèges dans un contexte microservices peut prendre plusieurs
formes : accès à un endpoint non autorisé, utilisation d'un token JWT avec des
claims forgés, ou accès à un service interne depuis un namespace non autorisé.

Les \textbf{AuthorizationPolicies} d'Istio permettent un contrôle d'accès
granulaire basé sur : les méthodes HTTP, les chemins URI, les namespaces source,
les ServiceAccounts SPIFFE et les claims JWT.

\medskip\noindent
\textbf{Distinction authentification / autorisation :} Un token JWT valide mais
ne contenant pas les claims requis sera accepté par \texttt{RequestAuthentication}
(validation de signature réussie) mais rejeté par l'\texttt{AuthorizationPolicy}
(code \texttt{403}). Cette séparation entre les deux niveaux de contrôle est
fondamentale pour une défense en profondeur.

% ─────────────────────────────────────────────────────────────────────────────
\section{Synthèse des Risques et Contre-mesures}
% ─────────────────────────────────────────────────────────────────────────────

Le tableau suivant synthétise les onze scénarios d'attaque analysés et
expérimentés, avec les mécanismes Istio responsables du blocage :

\begin{table}[H]
  \centering
  \caption{Scénarios d'attaque STRIDE et contre-mesures Istio}
  \label{tab:stride_scenarios}
  \renewcommand{\arraystretch}{1.4}
  \begin{tabularx}{\textwidth}{
    >{\bfseries}C{0.6cm} C{2cm} L{3.2cm} L{4cm} C{3cm} C{1.2cm}
  }
    \toprule
    \rowcolor{gray!15}
    ID & STRIDE & Scénario & Mécanisme Istio & Résultat & Code \\
    \midrule
    T1  & Elevation   & Accès \texttt{/charge} sans JWT
        & AuthorizationPolicy + RequestAuthentication
        & \textbf{BLOQUÉ} \checkmark & 403 \\
    \rowcolor{gray!10}
    T2  & Spoofing    & Pod sans sidecar $\rightarrow$ payment-api
        & mTLS STRICT PeerAuthentication
        & \textbf{BLOQUÉ} \checkmark & Reset \\
    T3  & Elevation   & Cross-namespace $\rightarrow$ fraud-detection
        & NetworkPolicy + AuthZ
        & \textbf{BLOQUÉ} \checkmark & Reset \\
    \rowcolor{gray!10}
    T4  & Spoofing    & Wrong SPIFFE SA $\rightarrow$ fraud-detection
        & SPIFFE principal validation
        & \textbf{BLOQUÉ} \checkmark & 403 \\
    T5  & Tampering   & JWT signature corrompue
        & RequestAuthentication sig. verify
        & \textbf{BLOQUÉ} \checkmark & 401 \\
    \rowcolor{gray!10}
    T6  & Tampering   & JWT issuer non reconnu
        & RequestAuthentication issuer check
        & \textbf{BLOQUÉ} \checkmark & 401 \\
    T7  & Repudiation & Traçabilité attaques
        & Envoy access log horodaté
        & \textbf{TRACÉ} \checkmark & --- \\
    \rowcolor{gray!10}
    T8  & Info Discl. & Interception trafic inter-pod
        & mTLS chiffrement TLS~1.3
        & \textbf{CHIFFRÉ} \checkmark & --- \\
    T9  & Info Discl. & SA token exfiltration
        & RBAC Kubernetes + SPIFFE
        & \textbf{LIMITÉ} \checkmark & 403 \\
    \rowcolor{gray!10}
    T10 & DoS         & Flood 200 requêtes/s
        & Connection pool DestinationRule
        & \textbf{DÉGRADÉ} $\triangle$ & 5xx \\
    T11 & DoS         & Rate limiting EnvoyFilter
        & EnvoyFilter token bucket 20/s
        & \textbf{PROTÉGÉ} \checkmark & 429 \\
    \bottomrule
  \end{tabularx}
\end{table}

% ─────────────────────────────────────────────────────────────────────────────
\begin{samepage}
\section*{Conclusion}
% ─────────────────────────────────────────────────────────────────────────────

L'analyse STRIDE révèle qu'Istio offre une couverture complète des six catégories
de menaces identifiées. La combinaison \textbf{mTLS + JWT + AuthorizationPolicy
+ SPIFFE + NetworkPolicy} crée une défense en profondeur où chaque couche
compense les limitations des autres. Sur les onze scénarios expérimentés, neuf
sont intégralement bloqués, un est tracé et un est partiellement atténué,
confirmant la robustesse de l'architecture déployée.
\end{samepage}